\documentclass[%
% jmp,
% bmf,
% sd,
% rsi,
twocolumn,
amsmath,
amssymb,
groupedaddress,
% preprint,%
reprint,%
% author-year,%
% author-numerical,%
% Conference Proceedings
]{revtex4-2}

\usepackage{bm}
\usepackage{listings}
\usepackage{subcaption}

\lstset{language=Python,keywordstyle={\bfseries \color{blue}}}

\usepackage{tikz}
\usetikzlibrary{arrows.meta,
	chains,
	matrix,
  decorations,
	positioning,
  calc,
  decorations.pathmorphing,
  shapes.geometric}

\definecolor{b0}{HTML}{6200EA}
\definecolor{b1}{HTML}{651FFF}
\definecolor{b2}{HTML}{7C4DFF}
\definecolor{b3}{HTML}{B388FF}
\definecolor{g0}{HTML}{B9F6CA}
\definecolor{g1}{HTML}{69F0AE}
\definecolor{g2}{HTML}{00E676}
\definecolor{g3}{HTML}{00C853}
\definecolor{r0}{HTML}{FF8A80}
\definecolor{r1}{HTML}{FF5252}
\definecolor{r2}{HTML}{FF1744}
\definecolor{r3}{HTML}{C51162}
\definecolor{lr0}{HTML}{FF9E80}
\definecolor{lr1}{HTML}{FF6E40}
\definecolor{lr2}{HTML}{FF3D00}
\definecolor{lr3}{HTML}{DD2C00}
\definecolor{y0}{HTML}{FFE57F}
\definecolor{y1}{HTML}{FFD740}
\definecolor{y2}{HTML}{FFC400}
\definecolor{y3}{HTML}{FFAB00}


\renewcommand{\d}{{\rm d}}
\newcommand{\w}{\omega}
\newcommand{\wti}{\widetilde}
\newcommand{\ti}{\tilde}
\newcommand{\B}{\mathrm{B}}
\newcommand{\D}{\mathrm{D}}
\newcommand{\Bs}{\mathrm{Bose}}
\newcommand{\s}{\mathrm{S}}
\newcommand{\tT}{\mathrm{T}}
\newcommand{\tB}{\mathrm{B}}
\newcommand{\tW}{\mathrm{W}}
\newcommand{\tS}{\mathrm{S}}
\newcommand{\st}{\mathrm{st}}
\newcommand{\te}{\mathrm{e}}
\newcommand{\tg}{\mathrm{g}}
\newcommand{\T}{\mathrm{T}}
\newcommand{\tR}{\mathrm{R}}
\newcommand{\tL}{\mathrm{L}}
\newcommand{\tM}{\mathrm{M}}
\newcommand{\tc}{\mathrm{c}}
\newcommand{\app}{\mathrm{app}}
\newcommand{\tot}{\mathrm{tot}}
\newcommand{\SB}{\mathrm{SB}}
\newcommand{\tr}{\mathrm{tr}}
\newcommand{\Tr}{\mathrm{Tr}}
\newcommand{\dg}{\dagger}
\newcommand{\la}{\langle}
\newcommand{\ra}{\rangle}
\renewcommand{\>}{\rangle}
\newcommand{\<}{\langle}
\newcommand{\lla}{\langle\langle}
\newcommand{\rra}{\rangle\rangle}
\newcommand{\La}{\big\la}
\newcommand{\Ra}{\big\ra}
\newcommand{\BO}{\mathrm{BO}}
\newcommand{\cm}{\mathrm{cm}^{\mathrm{-1}}}
% \newcommand{\Ref}[1]{Ref.\,\onlinecite{#1}}
\newcommand{\Sec}[1]{Sec.\,\ref{#1}}
\newcommand{\nl}{\nonumber \\}
\newcommand{\be}{\begin{equation}}
\newcommand{\ee}{\end{equation}}
\newcommand{\bsube}{\begin{subequations}}
\newcommand{\esube}{\end{subequations}}
\newcommand{\Eq}[1]{Equation \,(\ref{#1})}
\newcommand{\eq}[1]{Eq.\,(\ref{#1})}
\newcommand{\Eqs}[1]{Equations .\,(\ref{#1})}
\newcommand{\eqs}[1]{Equations .\,(\ref{#1})}
\newcommand{\Fig}[1]{Fig.\,\ref{#1}}
\newcommand{\Tab}[1]{Table \ref{#1}}
\newcommand{\obarplus}{\mbox{\tiny$\oplus$}}
\newcommand{\obarminus}{\mbox{\tiny$\ominus$}}
\newcommand{\greater}{\mathrm{$>$}}
\newcommand{\lesser}{\mathrm{$<$}}
\newcommand{\lgter}{\mathrm{$\lessgtr$}}
\newcommand{\gler}{\mathrm{$\gtrless$}}
\newcommand{\WF}{\mathrm{WF}}
\newcommand{\KS}{\mathrm{KS}}
\newcommand{\XC}{\mathrm{XC}}
\renewcommand{\L}{\mathrm{L}}
\renewcommand{\P}{\mathrm{P}}
\newcommand{\lyp}{\mathrm{b3lyp}}
\newcommand{\LDA}{\mathrm{LDA}}
\newcommand{\brm}[1]{\bm{\mathrm{#1}}}

\newcommand{\blue}[1]{\textcolor{blue}{#1}}
\newcommand{\red}[1]{\textcolor{red}{#1}}


\begin{document}


\title{
	A molecule
}
\author{XX}

\date{\today}

\begin{abstract}
	%%
	{
		TODO
	}
\end{abstract}

\maketitle

\allowdisplaybreaks

\section{Introduction}
The well-known Kohn--Sham equation \cite{Koh65A1133} is
\begin{align}
	\label{eq:KS}
	\Big( - \frac{1}{2} \nabla^{2} + v_{\mathrm{ext}} + v_{\mathrm{H}} + v_{\mathrm{XC}}\Big) \phi_{i} = \epsilon_{i} \phi_{i}.
\end{align}
% 
Where $\phi_{i}$ is the Kohn--Sham orbital, $\epsilon_{i}$ is the Kohn--Sham orbital energy, $v_{\mathrm{ext}}$ is the external potential (such as the electrostatic potential of nuclei), $v_{\mathrm{H}}$ is the electrostatic potential of the electron, and $v_{\mathrm{XC}}$ is the exchange-correlation potential.
% 
According to the Hohenberg-Kohn theorem \cite{Hoh64B864}, the external potential, $v_{\mathrm{ext}}$ is the unique functional of the ground-state electron density, $\rho = n_{i} \vert \phi_{i} \vert^{2}$, where $n_{i}$ is the occupation number of the Kohn-Sham orbital.
% 
So, once this unique functional is obtained, the Kohn-Sham equation can be solved iteratively to obtain the electron density and the other properties of the molecule.

One of the challenging tasks of the DFT calculation is to find the unique functional between the exact potential $v_{\mathrm{XC}}$ and the electron density $\rho$.
% 
In this article, we combine the machine learning (ML) method and density functional theory (DFT) to predict the properties of the molecules such as the energy, the dipole, the bond length, the bond dissociation energy, the band gap, and so on.

Among various methods to solve \eq{eq:KS}, the self-consistent field (SCF) method is widely used.
%
In the SCF method, the Kohn-Sham equation is solved iteratively utilizing the so-called potential term $v_{\mathrm{XC}} = \delta E_{\mathrm{XC}} / \delta \rho$, where $E_{\mathrm{XC}}$ is the exchange-correlation energy functional.
% 


\section{The NEURAL NETWORK}
We use the transformer \cite{Vas17A5998} as the neural network to learn the functional between the electron density and the energy density of the molecule.
% 
After the training, the neural network can predict the energy of the molecule given the electron density and then can be used to optimize the molecular structure.

\subsection{The architecture of the neural network}
\begin{figure*}
	\tikzset{%
		block/.style={rectangle, draw, fill=blue!20, text centered, rounded corners, minimum height=0.6cm},
		block1/.style={rectangle, fill=white, text centered, rounded corners, minimum height=0.6cm},
	}
	\begin{tikzpicture}[auto, node distance=1.5cm]
		% We start by placing the blocks
		\node [block] (input) {${\bm{e}_{\rm{cube}}(\brm{r}_i)}$};
		\node [block, right=of input] (transformer) {Transformer};
		\node [block, right=of transformer] (flatten) {Flatten};
		\node [block, right=of flatten] (linear) {Linear};
		\node [block, right=of linear, xshift=0.5cm] (output) {MLP};
		\node [block1, right=of output] (output_final) {};
		\node [block1, below=of input, yshift=0.5cm] (input1) {};
		\node [block, below=of transformer, yshift=0.5cm, xshift=1cm] (flatten1) {Extraction};
		\node [block, below=of flatten, yshift=0.5cm, xshift=1.5cm] (linear1) {Linear};
		\node [block1, below=of output, yshift=0.5cm] (output1) {};
		% connecting the blocks
		\draw [->] (input) -- node {$(N, 27, 4)$} (transformer);
		\draw [->] (transformer) -- node {$(N, 27, 4)$} (flatten);
		\draw [->] (flatten) -- node {$(N, 108)$} (linear);
		\draw [->] (linear) -- node {$(N, 1)$} (output);
		\draw [->] (output) -- node {$(N, 1)$} (output_final);
		\draw [->, rounded corners] (input.south) -- (input1.center) -- node {$(N, 27, 4)$} (flatten1);
		\draw [->] (flatten1) -- node {$(N, 4)$} (linear1);
		\draw [->, rounded corners] (linear1) -- node {$(N, 1)$} ([xshift=-0.75cm] output1.center) -- ([xshift=-0.75cm] output.center);
		% \draw [->] (embedding) -- node {$\bm{h}_0$} (transformer);
		% \draw [->] (transformer) -- node {$\bm{h}_L$} (linear);
		% \draw [->] (linear) -- node {$E(\brm{r}_i)$} (output);
	\end{tikzpicture}
	\caption{
		% 
		The architecture of the neural network.
		% 
		The transformer we use here is only composed of multilayer of the attention part of the original transformer \cite{Vas17A5998}.
		%
		The Extraction block extracts the energy density at the center (${\bm{e}(\brm{r}_i)}$) from the input.
		% 
		The MLP block is a multilayer perceptron with two hidden layers, before the MLP block the two inputs are concatenated.
		%
	}
	\label{fig:arch}
\end{figure*}

The architecture of the neural network is shown in \Fig{fig:arch}.
%
The input of the neural network is a set of the energy density ${\tilde{\bm{e}}_{\rm{cube}}(\brm{r}_i)}$ with the size of $(N, 27, 4)$, and the output is the energy density ${e(\brm{r}_i)}$ with the size of $(N, 1)$.
% 
N is the number of the grid points in the normal dft procedure.
%
\begin{align}
	\mathbf{r}_{i} =
	 & (x_{i},y_{i},z_{i})                                                                              \\
	\bm{e}_{\rm{cube}}(\brm{r}_{i}) =
	 & \big\{ \bm{e}(\mathbf{r}'_{i}) \text{ for } \mathbf{r}'_{i} \in \mathbf{r}_{i;\rm{cube}} \big\}. \\
	\bm{e}(\brm{r}) =
	 & (e_{\rm{LDA}}(\brm{r}), e_{\rm{VWN3}}(\brm{r}), e_{\rm{B88}}(\brm{r}), e_{\rm{LYP}}(\brm{r})),   \\
	\mathbf{r}_{i;\rm{cube}} =
	 & \big\{ \mathbf{r}'_{i} = (x',y',z') \vert \nl
	 & x' \in [x_{i}-d,x_{i},x_{i}+d], \nl
	 & y' \in [y_{i}-d,y_{i},y_{i}+d], \nl
	 & z' \in [z_{i}-d,z_{i},z_{i}+d] \big\}.
\end{align}

\subsection{DATASET AND THE LOSS FUNCTION}

The loss function of $E_{\mathrm{target}}$ is defined as
\begin{align}
	\label{eq:loss}
	L_{e} =
	 & \sqrt{N_{\rm{atom}}} \Big| \int \d \brm{r}_{i} \big(e_{\mathrm{target}}(\brm{r}_{i}) - e_{\mathrm{predict}}(\brm{r}_{i})\big) \Big| \nl
	 & + \lambda_{\rm{ABS}} \sqrt{N_{\rm{atom}}} \int \d \brm{r}_{i} \Big| e_{\mathrm{target}}(\brm{r}_{i}) - e_{\mathrm{predict}}(\brm{r}_{i}) \Big| \nl
	 & + \lambda_{\rm{AE}} \sqrt{N_{\rm{atom}}} \Big| E^{\rm{AE}}_{\mathrm{target}}(\brm{r}_{i}) - E^{\rm{AE}}_{\mathrm{predict}}(\brm{r}_{i}) \Big| .
\end{align}
where $N_{\rm{atom}}$ is the number of atoms in the molecule and $E^{\rm{AE}}_{\mathrm{predict}}$ and $E^{\rm{AE}}_{\mathrm{target}}$ is the atomic energy of the predicted and target values, respectively.
%

\subsection{Result}

\bibliography{./ref}

\pagebreak

\appendix

\section{training details}
\subsection{dataset}

The dataset we use are listed in \Tab{tab:dataset}.
\begin{table}
	\caption{
		The dataset we use.
		%
		\label{tab:dataset}
	}
	\begin{ruledtabular}
		\begin{tabular}{lcccccc}
			\textrm{Dataset} & \textrm{Size} & \textrm{Heavy Atoms} \\
			\colrule
			W4-11            & 100           & $\leq 2$             \\
		\end{tabular}
	\end{ruledtabular}
\end{table}

\end{document}
